\section{Tractable Special Cases: When You Can Solve It}\label{sec:tractable}

Formally: this section gives sufficient structure/representation conditions for polynomial-time SUFFICIENCY-CHECK.

We distinguish the encodings of Section~\ref{sec:encoding}. The tractability results below state the model assumption explicitly.
Structural insight: hardness dissolves exactly when the full decision boundary $s \mapsto \Opt(s)$ is recoverable in polynomial time from the input representation; the six cases below instantiate this single principle.
Concretely, each tractable regime corresponds to a specific structural insight (explicit boundary exposure, additive separability, tensor factorization, tree structure, bounded treewidth, or coordinate symmetry) that removes the hardness witness; this supports reading the general-case hardness as missing structural access in the current representation rather than as an intrinsic semantic barrier.
In physical terms, these are the regimes where required structure is directly observable or factorized so that verification work scales with exposed local structure rather than hidden global combinations on encoded physical instances.\claimstamp{T\ref{thm:tractable},\ref{thm:physical-bridge-bundle};P\ref{prop:physical-claim-transport}}{E,TA,AR}
\leanmeta{\LHrng{CC}{69}{73}, \LH{CT4}, \LH{CT6}.}

Informally: these cases are tractable because useful structure is visible.

\begin{theorem}[Tractable Subcases]\label{thm:tractable}
SUFFICIENCY-CHECK is polynomial-time solvable in the following cases:
\begin{enumerate}
\item \textbf{Bounded actions (explicit-state encoding):} The input contains the full utility table over $A \times S$. SUFFICIENCY-CHECK runs in $O(|S|^2|A|)$ time; if $|A|$ is constant, $O(|S|^2)$.
\item \textbf{Separable utility (rank 1):} $U(a, s) = f(a) + g(s)$. The optimal action is state-independent, so $I = \emptyset$ is always sufficient. Complexity: $O(1)$ for sufficiency verification.
\item \textbf{Low tensor rank:} $U(a,s) = \sum_{r=1}^R w_r \cdot f_r(a) \cdot \prod_{i=1}^n g_{ri}(s_i)$ for bounded $R$. Complexity: $O(|A| \cdot R \cdot n)$.
\item \textbf{Tree-structured dependencies:} Dependencies form a rooted tree:
\[
U(a,s) = \sum_i u_i\bigl(a, s_i, s_{\mathrm{parent}(i)}\bigr),
\]
with the root term depending only on $(a, s_{\mathrm{root}})$ and all $u_i$ given explicitly. Complexity: $O(n \cdot |A| \cdot k_{\max})$.
\item \textbf{Bounded treewidth:} Utility decomposes as pairwise interactions on an interaction graph of treewidth $w$. Complexity: $O(n \cdot k^{w+1})$ via CSP algorithms~\cite{bodlaender1993tourist, freuder1990complexity}.
\item \textbf{Coordinate symmetry:} State space is $S = [k]^d$ with utility invariant under coordinate permutations. The effective state space reduces from $k^d$ to $\binom{d+k-1}{k-1}$ orbit types. Complexity: $O\bigl(\binom{d+k-1}{k-1}^2\bigr)$, polynomial in $d$ for fixed $k$.
\end{enumerate}
\leanmeta{\LHrng{CC}{70}{73}.}
\end{theorem}

\subsection{Bounded Actions (Explicit-State)}
\leanmeta{\LH{CC70}.}

\begin{proof}[Proof of Part 1]
Given the full table of $U(a,s)$, compute $\Opt(s)$ for all $s \in S$ in $O(|S||A|)$ time. For SUFFICIENCY-CHECK on a given $I$, verify that for all pairs $(s,s')$ with $s_I = s'_I$, we have $\Opt(s) = \Opt(s')$. This takes $O(|S|^2|A|)$ time by direct enumeration and is polynomial in the explicit input length. If $|A|$ is constant, the runtime is $O(|S|^2)$.
\end{proof}

\subsection{Separable Utility (Rank 1)}
\leanmeta{\LH{CC71}.}

\begin{proof}[Proof of Part 2]
If $U(a, s) = f(a) + g(s)$, then $\Opt(s) = \arg\max_{a \in A} f(a)$. The optimal action is independent of the state. Thus $I = \emptyset$ is always sufficient, and sufficiency verification is $O(1)$.
\end{proof}

\subsection{Low Tensor Rank}
\leanmeta{\LH{CC71}.}

\begin{proof}[Proof of Part 3]
When utility has tensor rank $R$: $U(a,s) = \sum_{r=1}^R w_r f_r(a) \prod_{i=1}^n g_{ri}(s_i)$, the problem reduces to factored tensor computation. For each action, the optimal value computation requires $O(R \cdot n)$ operations. Total complexity: $O(|A| \cdot R \cdot n)$, polynomial for bounded $R$. This reduction cites standard tensor contraction algorithms~\cite{kolda2009tensor, cichocki2016tensor}.
\end{proof}

\subsection{Tree-Structured Dependencies}
\leanmeta{\LH{CC73}.}

\begin{proof}[Proof of Part 4]
Assume the tree decomposition and explicit local tables as stated. For each node $i$ and each value of its parent coordinate, compute the set of actions that are optimal for some assignment of the subtree rooted at $i$. This is a bottom-up dynamic program that combines local tables with child summaries; each table lookup is explicit in the input. A coordinate is relevant if and only if varying its value changes the resulting optimal action set. The total runtime is polynomial in $n$, $|A|$, and the size of the local tables.
\end{proof}

\subsection{Bounded Treewidth}
\leanmeta{\LH{CC73}, \LH{CC75}.}

\begin{proof}[Proof of Part 5]
When utility decomposes as pairwise interactions $U(a,s) = \sum_{(i,j)} u_{ij}(a, s_i, s_j)$, the interaction graph has an edge between coordinates $i,j$ iff $u_{ij}$ is nontrivial. If this graph has treewidth $w$, SUFFICIENCY-CHECK reduces to a constraint satisfaction problem on a bounded-treewidth graph. By standard CSP algorithms~\cite{bodlaender1993tourist, freuder1990complexity}, this is solvable in $O(n \cdot k^{w+1})$ time.
\end{proof}

\subsection{Coordinate Symmetry}
\leanmeta{\LH{CC72}.}

\begin{proof}[Proof of Part 6]
For state space $S = [k]^d$ with coordinate-permutation-invariant utility, two states have the same optimal actions iff they have the same \emph{orbit type} (value histogram). The number of orbit types is $\binom{d+k-1}{k-1}$ by stars-and-bars. Sufficiency checking reduces to comparing optimal actions across orbit type pairs. Total complexity: $O\bigl(\binom{d+k-1}{k-1}^2 \cdot |A|\bigr)$, polynomial in $d$ for fixed $k$.
\end{proof}

\subsection{Practical Implications}

Formally: this subsection states detect-then-check closure for the declared tractable classes.

\begin{center}
\begin{tabular}{ll}
\toprule
\textbf{Condition} & \textbf{Examples} \\
\midrule
Bounded actions & Small or fully enumerated state spaces \\
Separable utility & Additive costs, linear models \\
Low tensor rank & Factored decision trees, sparse interactions \\
Tree-structured & Hierarchies, causal trees, DAGs \\
Bounded treewidth & Local dependency graphs, grid-like structures \\
Coordinate symmetry & Exchangeable particles, counting problems \\
\bottomrule
\end{tabular}
\end{center}

\begin{proposition}[Heuristic Reusability]\label{prop:heuristic-reusability}
\claimstamp{P\ref{prop:heuristic-reusability}}{H=tractable-structured}
Let $\mathcal{C}$ be a structure class for which SUFFICIENCY-CHECK is polynomial (Theorem~\ref{thm:tractable}). If membership in $\mathcal{C}$ is decidable in polynomial time, then the combined detect-then-check procedure is a sound, polynomial-time heuristic applicable to all future instances in $\mathcal{C}$.
For each subcase of Theorem~\ref{thm:tractable}, structure detection is polynomial (under the declared representation assumptions).
\leanmeta{\LH{CC7}, \LH{CC53}, \LH{CC55}, \LH{CC74}.}
\end{proposition}

\begin{remark}[Heuristics and Integrity]
Proposition~\ref{prop:heuristic-reusability} removes the complexity-theoretic obstruction to integrity-preserving action on detectable tractable instances: integrity no longer \emph{forces} abstention. Whether a declared controller executes the action still requires competence (Definition~\ref{def:competence-regime}): budget and coverage must also be satisfied. A controller that detects structure class $\mathcal{C}$, applies the corresponding polynomial checker, and abstains when detection fails maintains integrity; it never claims sufficiency without verification. Mistaking the heuristic for the general solution, claiming polynomial-time competence on a coNP-hard problem, violates integrity (Proposition~\ref{prop:integrity-resource-bound}).
\end{remark}

Informally: reusable heuristics are reliable only after the stated structure test; otherwise hard-case limits still apply.
